\documentclass[12pt, letterpaper]{article}

% --- PREÁMBULO ---
\usepackage[utf8]{inputenc}
\usepackage[T1]{fontenc}
\usepackage[spanish]{babel}
% Márgenes de 1 pulgada en todos los lados
\usepackage[letterpaper, margin=1in]{geometry}
\usepackage{graphicx}
\usepackage[most]{tcolorbox}
\usepackage{enumitem}
\usepackage{amssymb}
\usepackage{amsmath}
\usepackage{multicol}

% --- PAQUETE PARA FONDO ---
\usepackage{eso-pic}



% --- FUENTE ---
\usepackage{helvet}
\renewcommand{\familydefault}{\sfdefault}

% --- COLORES INSTITUCIONALES ---
\definecolor{valdark}{RGB}{50, 50, 50}

% --- COMANDO PARA CASILLAS DE RESPUESTA PEQUEÑAS ---
\newcommand{\boxfill}[1][1.5cm]{\tcbox[on line, size=minimal, colframe=valdark!80, colback=white, boxrule=0.5pt, arc=1mm, boxsep=3pt, left=2pt, right=2pt, top=2pt, bottom=2pt]{\hspace*{#1}\vphantom{M}}}

\begin{document}
	
	% --- ENCABEZADO FORMATO OFICIAL TALLER DE REPASO ---
	\begin{center}
		\textbf{TALLER DE REPASO}\\[0.1cm]
		Matemáticas - Unidad 7: Adición y Sustracción de Fracciones
	\end{center}
	
	
	% --- PUNTO 1 ---
	\noindent \textbf{1. Resuelve las siguientes sumas y restas.}
	
	\vspace{0.4cm}
	\begin{multicols}{3}
		\begin{enumerate}[label=\textbf{\Alph*)}, itemsep=0.8cm]
			\item $\dfrac{4}{9} + \dfrac{3}{9} = \dfrac{\boxfill[0.8cm]}{\boxfill[0.8cm]}$
			
			\item $\dfrac{8}{12} - \dfrac{5}{12} = \dfrac{\boxfill[0.8cm]}{\boxfill[0.8cm]}$
			
			\item $\dfrac{2}{7} + \dfrac{4}{7} = \dfrac{\boxfill[0.8cm]}{\boxfill[0.8cm]}$
		\end{enumerate}
	\end{multicols}
	
	\vspace{0.8cm}
	
	% --- PUNTO 2 ---
	\noindent \textbf{2. Resuelve las siguientes sumas}
	
	\vspace{0.5cm}
	\begin{multicols}{2}
		\begin{enumerate}[label=\textbf{\Alph*)}, itemsep=1.5cm]
			\item $\dfrac{1}{3} + \dfrac{1}{4} = \dfrac{\boxfill[0.8cm] + \boxfill[0.8cm]}{\boxfill[1cm]} = \dfrac{\boxfill[1cm]}{\boxfill[1cm]}$
			
			\item $\dfrac{2}{5} + \dfrac{1}{2} = \dfrac{\boxfill[0.8cm] + \boxfill[0.8cm]}{\boxfill[1cm]} = \dfrac{\boxfill[1cm]}{\boxfill[1cm]}$
		\end{enumerate}
	\end{multicols}
	
	\vspace{0.8cm}
	
	% --- PUNTO 3 ---
	\noindent \textbf{3. Resuelve las siguientes restas}
	
	\vspace{0.5cm}
	\begin{multicols}{2}
		\begin{enumerate}[label=\textbf{\Alph*)}, itemsep=1.5cm]
			\item $\dfrac{4}{5} - \dfrac{1}{3} = \dfrac{\boxfill[0.8cm] - \boxfill[0.8cm]}{\boxfill[1cm]} = \dfrac{\boxfill[1cm]}{\boxfill[1cm]}$
			
			\item $\dfrac{3}{4} - \dfrac{1}{2} = \dfrac{\boxfill[0.8cm] - \boxfill[0.8cm]}{\boxfill[1cm]} = \dfrac{\boxfill[1cm]}{\boxfill[1cm]}$
		\end{enumerate}
	\end{multicols}
	
	\newpage
	
	% --- PUNTO 4 ---
	\noindent \textbf{4. Lee la siguiente situación, completa la operación y calcula el total. }
	
	\vspace{0.4cm}
	\noindent \textit{Para hacer una receta deliciosa, necesitas mezclar $\mathbf{\dfrac{2}{8}}$ de taza de azúcar blanca y $\mathbf{\dfrac{3}{8}}$ de taza de azúcar morena. ¿Cuánta azúcar usaste en total para la mezcla?}
	
	\vspace{0.6cm}
	\begin{center}
		\begin{tabular}{c c c}
			\textbf{Operación:} & \hspace{2cm} & \textbf{Respuesta:} \\[0.6cm]
			\Large $\dfrac{\boxfill[0.8cm]}{\boxfill[0.8cm]} + \dfrac{\boxfill[0.8cm]}{\boxfill[0.8cm]} = \dfrac{\boxfill[0.8cm]}{\boxfill[0.8cm]}$ & & 
			Usaste $\dfrac{\boxfill[0.8cm]}{\boxfill[0.8cm]}$ de taza en total.
		\end{tabular}
	\end{center}
	
	\vspace{1.5cm}
	
	% --- PUNTO 5 ---
	\noindent \textbf{5. Lee la siguiente situación, completa la operación y calcula lo que sobra. }
	
	\vspace{0.4cm}
	\noindent \textit{Un galón de pintura estaba lleno hasta los $\mathbf{\dfrac{9}{10}}$ de su capacidad. Si usaste $\mathbf{\dfrac{3}{10}}$ del galón para pintar tu cuarto, ¿cuánta pintura te queda guardada en el galón?}
	
	\vspace{0.6cm}
	\begin{center}
		\begin{tabular}{c c c}
			\textbf{Operación:} & \hspace{2cm} & \textbf{Respuesta:} \\[0.6cm]
			\Large $\dfrac{\boxfill[0.8cm]}{\boxfill[0.8cm]} - \dfrac{\boxfill[0.8cm]}{\boxfill[0.8cm]} = \dfrac{\boxfill[0.8cm]}{\boxfill[0.8cm]}$ & & 
			Quedan $\dfrac{\boxfill[0.8cm]}{\boxfill[0.8cm]}$ del galón.
		\end{tabular}
	\end{center}
	
	
	
\end{document}